% Section 6.3, from lectures/L06/appendix/b_exercises.md.

\section{Exercises}
\label{c6:sec:exercises}\label{c6:app:b}

\begin{aside}
\textbf{How to work these.} These exercises build the driver's three contracts from
\secref{c6:sec:driver_stack}: the register map, the transport interface, and the driver interface.
All three are header-only declarations, so there is nothing to run yet; \chapref{7} implements the
driver over them and \chapref{8} brings the first tests. Exercise~\ref{c6:ex:4.1} then implements
the driver interface once, as a UART stub.

\textbf{How to check your work.} For now the check is simply that each header compiles. Nothing in
\file{fw/Makefile} compiles a header on its own yet (\code{make build} waits for the
\file{source/main.cpp} of \chapref{8}, and \code{make test} for the headers of \chapref{8}), so
compile each one by hand, from \file{fw/}:

\begin{shell}
cd fw
for h in driver/uart/register_map.hpp driver/transport/interface.hpp driver/uart/stub.hpp; do
    echo "#include \"$h\"" | g++ -std=c++17 -Wall -Wextra -Werror -Iinclude -fsyntax-only -x c++ -
done
\end{shell}

\noindent Every exercise in this chapter is code, so Appendix~\ref{sol:ch} has nothing to add here.
The names below are the contract the host suite of \chapref{8},
\file{fw/test/uart/uart_test.cpp}, binds to, and that suite is where these headers are first put to
work.
\end{aside}

Work in \file{fw/}. Create the following files:

\begin{console}
include/
    driver/
        uart/
            register_map.hpp
            interface.hpp
            stub.hpp
        transport/
            interface.hpp
\end{console}

The register map and the driver interface live in the \code{driver::uart} namespace; the transport
interface lives in \code{driver::transport}. The names below are the contract the tests of
\chapref{8} bind to, so keep them as written.

\textbf{Write it AVR-portable.} These headers compile on the host (g++) now and, from \chapref{9},
on the ATmega328P. \textbf{We deliberately skip several modern C++ features here because the AVR
toolchain does not support them}, and the same driver code has to build for both targets, so stay on
the subset that toolchain accepts. The reference point is the AVR/GNU C++ compiler bundled with
Microchip Studio, which is several GCC releases behind a desktop g++; something compiling on your
host, or even on a Linux avr-g++, does not mean it will compile there.

Include \code{<stdint.h>}, and \code{<stddef.h>} if you need \code{size_t}, and use the bare types
\code{uint8_t}, \code{uint16_t}, \code{uint32_t} and \code{size_t}: the AVR toolchain ships no C++
standard library, so there is no \code{<cstdint>} and no \code{std::} namespace at all, and anything
reaching for either will fail.

Write namespaces as nested \textbf{blocks}, \code{namespace driver { namespace uart { ... } }},
because the compact C++17 form \code{namespace driver::uart { ... }} does not compile there. Leave
out \code{[[nodiscard]]}, which it rejects for the same reason. What you can rely on is
\code{noexcept}, \code{override}, \code{final}, \code{= default}, and pure \code{= 0}.

% ------------------------------------------------------------------------------------------
\exerciseset{The Register Map}
\label{c6:set:1}

\exercise{\code{register_map.hpp}}{C++}
\label{c6:ex:1.1}
In \file{include/driver/uart/register_map.hpp}, transcribe the register map from Part~2 of the
protocol (\secref{proto:part2}) as \code{constexpr} constants. Plain \code{constexpr}, not
\code{inline constexpr}: inline \emph{variables} are a C++17 feature the AVR toolchain does not
have. At namespace scope \code{constexpr} already implies internal linkage, so each translation unit
gets its own copy of a compile-time constant, which is exactly what you want and costs nothing.
These are the same values \file{uart_def.vhd} holds on the other side of the wire; a mismatch is a
bug you would only find on the bench.

\task{a) Register indices}
In a nested namespace \code{reg}, add one \code{uint8_t} constant per register index, each one the
register's offset divided by 4: \code{STATUS = 0}, \code{CTRL = 1}, \code{BAUD_DIV = 2},
\code{TX_DATA = 3}, \code{RX_DATA = 4}, \code{RX_POP = 5} and \code{ERROR_FLAGS = 6}.

\task{b) \code{STATUS} bit positions}
In a nested namespace \code{status}, add one \code{uint8_t} \textbf{bit-position} constant per
\code{STATUS} bit, meaning the bit's index rather than a mask: \code{TX_READY = 0},
\code{RX_VALID = 1}, \code{ERROR = 2} and \code{TX_IDLE = 3}. These are the same positions
\file{uart_def.vhd} defines, so the two sides agree.

You form a mask by shifting at the use site: \code{1U << status::TX_READY} tests that bit, and
\code{status & (1U << status::RX_VALID)} is non-zero when RX has data.

\task{c) \code{CTRL} and error bit positions}
In a nested namespace \code{ctrl}, add \code{ENABLE = 0} as a \code{uint8_t} bit position. In a
nested namespace \code{error}, add \code{FRAMING = 0}, \code{PARITY = 1} and \code{OVERRUN = 2}, all
\code{uint8_t} bit positions. The VHDL package also defines parity, stop-bit and interrupt-mask
positions in \code{CTRL}; this driver does not use them yet, so it does not declare them.

Positions rather than masks keep the C++ constants identical to the VHDL \file{uart_def.vhd}, which
also stores positions, and the one-line \code{1U << pos} at each use site is explicit about which bit
you mean.

% ------------------------------------------------------------------------------------------
\exerciseset{The Transport Interface}
\label{c6:set:2}

\exercise{\code{driver::transport::Interface}}{C++}
\label{c6:ex:2.1}
In \file{include/driver/transport/interface.hpp}, design an abstract class
\code{driver::transport::Interface} representing a byte-level SPI transport. It exchanges raw bytes
over one transaction and knows nothing about registers. All methods except the destructor shall be
pure virtual (\code{= 0}).

\task{a) Destructor}
Add a destructor that is virtual, marked \code{noexcept}, and implemented using \code{= default}.

\task{b) Begin a transaction}
Add a pure virtual method \code{begin()} that starts a transaction by pulling \code{SS} low. It takes
no parameters, returns nothing, and cannot throw exceptions.

\task{c) Exchange one byte}
Add a pure virtual method \code{transfer()} that exchanges a single byte in full duplex. It takes one
parameter of type \code{uint8_t}, the byte to send, returns the \code{uint8_t} received at the same
time, and cannot throw exceptions.

\task{d) End a transaction}
Add a pure virtual method \code{end()} that ends the transaction by releasing \code{SS}. It takes no
parameters, returns nothing, and cannot throw exceptions.

% ------------------------------------------------------------------------------------------
\exerciseset{The Driver Interface}
\label{c6:set:3}

\exercise{\code{Interface}}{C++}
\label{c6:ex:3.1}
In \file{include/driver/uart/interface.hpp}, design an abstract class \code{driver::uart::Interface},
the driver's public API. All methods except the destructor shall be pure virtual.

\task{a) Destructor}
Add a destructor that is virtual, \code{noexcept}, and \code{= default}.

\task{b) Configure}
Add a pure virtual method \code{configure()} that sets the baud divider and enables the peripheral.
It takes one parameter of type \code{uint16_t}, the baud divider, returns nothing, and cannot throw
exceptions.

\task{c) Write a byte}
Add a pure virtual method \code{write()} that attempts to send one byte. It takes one parameter of
type \code{uint8_t}, returns \code{true} if the byte was accepted or \code{false} if the transmit
FIFO was full, and cannot throw exceptions.

\task{d) Read a byte}
Add a pure virtual method \code{read()} that attempts to receive one byte. It takes a reference to a
\code{uint8_t} where the received byte will be stored, returns \code{true} if a byte was received and
\code{false} otherwise, and cannot throw exceptions.

\task{e) Status and errors}
Add three more pure virtual methods, none of which throws. \code{status()} returns the raw
\code{STATUS} register as a \code{uint32_t}, takes no parameters, and does not modify the object
(\code{const}). \code{errorFlags()} returns the raw \code{ERROR_FLAGS} register as a
\code{uint32_t}, also taking no parameters and also \code{const}. \code{clearErrors()} clears the
error flags, takes no parameters, and returns nothing.

% ------------------------------------------------------------------------------------------
\exerciseset{The UART Stub}
\label{c6:set:4}

\exercise{\code{driver::uart::Stub}}{C++}
\label{c6:ex:4.1}
The three headers above are all declarations. This one is the first concrete class of the C++
half, and it is worth writing now rather than later, because it needs nothing except the interface
you just declared: a \textbf{UART stub}, a \code{driver::uart::Interface} a test can drive.

It plays back scripted bytes on \code{read()} and records the bytes an application \code{write()}s,
so a test can queue input, run something over it, and check what came out. It is the same idea as
the \code{driver::transport::Stub} you write in \chapref{8}, one layer up: that one fakes the wire
under the driver, this one fakes the driver under an application. Write it in
\file{include/driver/uart/stub.hpp}, header-only, in the \code{driver::uart} namespace, in the same
AVR-portable style as everything else here; it holds no dynamic memory, so it compiles for the
target as well as the host.

Because a UART is a byte stream, both directions are \textbf{FIFO}: bytes are received and recorded
in the order they arrive. But a test queues a bounded number of bytes and reads them once, so the
storage is two plain \textbf{linear buffers} with no wraparound to reason about: two fixed
\code{uint8_t} arrays, each with a length. The \textbf{RX} buffer is read front to back through an
advancing index, so \code{read()} hands back the byte at the index and advances it, and the length
is how many bytes were queued. The \textbf{TX} buffer is a write-only append log, so \code{write()}
appends one byte and bumps the length, and the test reads the whole log back afterwards.

\textbf{Member variables}, using a fixed capacity \code{OurBufLen} and \code{uint8_t} throughout,
with no \code{size_t}. \code{myRxBuf[OurBufLen]}, with \code{myRxLen} and \code{myRxIdx}, holds the
scripted bytes to receive and tracks how far \code{read()} has advanced through them.
\code{myTxBuf[OurBufLen]}, with \code{myTxLen}, holds the bytes the application has sent, appended
in order.

The third is \code{volatile bool& myStop}, a reference to a caller-owned flag, set \code{true} when
\code{read()} reaches the end of the RX buffer. That one needs a word of explanation this early: in
\chapref{10} you write an application whose \code{run(const volatile bool& stop)} loops until the flag is
set, and a single-threaded test has no other way to end it. The stub feeds every scripted byte and
then, once the input is exhausted, asks the loop to stop. The flag is \code{volatile} because it is
written outside the loop that reads it; \chapref{10} says why that matters. Because it is a reference member, copy and
move are deleted, as is the default constructor, since you need a stop flag to build one.

\textbf{Methods}, the six from \code{driver::uart::Interface} plus scripting helpers. The
constructor, \code{Stub(volatile bool& stop)}, is \code{explicit}, stores the stop reference, and leaves both
buffers empty. \code{configure(uint16_t)} does nothing, since the stub has no baud rate to set.
\code{write(uint8_t byte)} appends \code{byte} to \code{myTxBuf} if there is room and returns
\code{true}, as a UART with room in its TX FIFO would, and returns \code{false} if the buffer is
full. \code{read(uint8_t& byte)} hands back \code{myRxBuf[myRxIdx++]} and returns \code{true} when
\code{myRxIdx < myRxLen}, and otherwise sets \code{myStop} to \code{true} and returns \code{false};
that empty-case flag is the loop's off switch. \code{status()}, \code{errorFlags()} and
\code{clearErrors()} are trivial, returning \code{0} or doing nothing, because the applications
tested this way drive the UART through \code{read()} and \code{write()} rather than the status
register.

Three helpers exist for the tests. \code{injectRxByte(uint8_t byte)} appends one scripted byte to
\code{myRxBuf}, returning \code{false} if it is full, and is called once per byte the application
should receive. \code{txBuf()} and \code{txLen()} return a pointer to the recorded TX bytes and how
many there are, so a test reads \code{txBuf()[i]} for \code{i < txLen()} in the order they were
sent. And \code{reset()} zeroes both lengths, the read index and the stop flag, so the stub can be
reused across tests.

Nothing exercises it yet, which is the honest state of a test double written before the thing it
doubles for. It compiles, and in \chapref{10} it becomes the entire test harness for
\code{app::EchoNode}.

\subsection*{What's ahead}
In \chapref{7} you implement \code{driver::uart::Uart}, which inherits this interface and drives it
through the register protocol over a \code{driver::transport::Interface}; \chapref{8} then writes the
scripted \code{driver::transport::Stub} and gets the host tests that check it green. The UART stub
written here waits until \chapref{10}, where it becomes the test harness for \code{app::EchoNode}.
